\documentclass[reqno]{amsart}
\usepackage[T1]{fontenc}
\usepackage[textwidth=6in,textheight=8.35in,centering]{geometry}
\usepackage{amsmath,amssymb,amsthm}
\usepackage{xr-hyper}
\usepackage[hidelinks]{hyperref}
\externaldocument[ps-]{20261125T010500--poincare-series-and-near-orthogonality__workbook}
\externaldocument[ob-]{20261126T040100--spaces-forms-and-twists__workbook}
\externaldocument[m1-]{20261126T051900--unfolding-twice-gram-operators-as-sums-over-rational-points__workbook}
\externaldocument[b1-]{20261126T081500--a-bessel-function-at-a-real-place__workbook}
\externaldocument[bp-]{20261126T020210--bessel-poincare-series-and-their-gram-operators__workbook}
\externaldocument[c2-]{20261126T050500--equal-multiplicities__workbook}

\newcounter{chap}
% Chapter numbers, generated from 20261125T012000--grh-workbook__guide.org
% by scripts/make-chapters.py.  Do not edit: change the order of the
% index and rerun the script (./preview-site.sh and CI do this).
%
% Usage in a sheet, after \newcounter{chap}:
%   % Chapter numbers, generated from 20261125T012000--grh-workbook__guide.org
% by scripts/make-chapters.py.  Do not edit: change the order of the
% index and rerun the script (./preview-site.sh and CI do this).
%
% Usage in a sheet, after \newcounter{chap}:
%   % Chapter numbers, generated from 20261125T012000--grh-workbook__guide.org
% by scripts/make-chapters.py.  Do not edit: change the order of the
% index and rerun the script (./preview-site.sh and CI do this).
%
% Usage in a sheet, after \newcounter{chap}:
%   \input{chapters}\chapnumber{<sheet ID>}
\makeatletter
\providecommand\chapters@set[2]{\@namedef{chapters@#1}{#2}}
\chapters@set{20261125T001000}{1}% zeros-to-ramp-sums
\chapters@set{20261125T001100}{2}% inert-primes-and-estimate-b
\chapters@set{20261125T005000}{3}% galois-fourier-and-averaging
\chapters@set{20261125T005500}{4}% columns-from-a-sphere
\chapters@set{20261125T000000}{5}% slope-method
\chapters@set{20261126T030100}{6}% a-twist-for-each-prime
\chapters@set{20261125T003000}{7}% the-chain-and-the-square-root
\chapters@set{20261125T010000}{8}% rank-lemma-and-zero-estimate
\chapters@set{20261125T010500}{9}% poincare-series-and-near-orthogonality
\chapters@set{20261124T223000}{10}% elliptic-curves-for-the-proof
\chapters@set{20261125T009500}{11}% fock-model-and-theta-intertwiners
\chapters@set{20261124T223100}{12}% complex-multiplication-and-inert-reduction
\chapters@set{20261125T008000}{13}% one-cm-determinant
\chapters@set{20261124T214500}{14}% supersingular-theta-divisibility
\chapters@set{20261124T214600}{15}% the-isotropic-plane
\chapters@set{20261125T013700}{16}% pairings-on-a-formal-group-of-height-two
\chapters@set{20261125T013500}{17}% hecke-nondegeneracy-on-the-tree
\chapters@set{20261126T030200}{18}% the-bound-is-attained
\chapters@set{20261125T009000}{19}% kronecker-limit-and-chowla-selberg
\chapters@set{20261125T002000}{20}% definite-theta-lifts
\chapters@set{20261125T015500}{21}% siegel-weil-in-the-smallest-case
\chapters@set{20261125T002100}{22}% definite-theta-lifts-and-norms
\chapters@set{20261125T016500}{23}% hecke-operators-and-theta-lifts
\chapters@set{20261125T016700}{24}% divisibility-along-the-chain
\chapters@set{20261125T004000}{25}% denominators-and-sha
\chapters@set{20261125T004100}{26}% tate-twists-and-sha
\chapters@set{20261126T020430}{27}% the-minus-part-by-gauss-sums
\chapters@set{20261125T014100}{28}% congruence-modules
\chapters@set{20261125T014500}{29}% eisenstein-congruences-and-selmer-classes
\chapters@set{20261125T014700}{30}% pseudocharacters-and-the-length-bound
\chapters@set{20261125T011000}{31}% galois-representations-black-box
\chapters@set{20261126T030300}{32}% the-data
\chapters@set{20261126T020210}{33}% bessel-poincare-series-and-their-gram-operators
\chapters@set{20261126T030400}{34}% the-seesaw-identity
\chapters@set{20261126T030500}{35}% divisibility-of-the-composite-lift
\chapters@set{20261126T030600}{36}% theorem-a-short-forms
\chapters@set{20261126T030700}{37}% the-determinant-of-a-chain
\chapters@set{20261126T030800}{38}% the-cost-of-one-lift
\chapters@set{20261126T030900}{39}% denominators-of-projections-to-lifts
\chapters@set{20261126T031000}{40}% twists-with-few-congruences
\chapters@set{20261126T031100}{41}% theorem-c-and-the-main-inequality
\providecommand\chapnumber[1]{%
  \@ifundefined{chapters@#1}%
    {\PackageError{chapters}{Sheet #1 is not listed in the index}%
      {List it in 20261125T012000--grh-workbook__guide.org\MessageBreak
       and run scripts/make-chapters.py.}}%
    {\edef\thechap{\@nameuse{chapters@#1}}}}
\makeatother
\chapnumber{<sheet ID>}
\makeatletter
\providecommand\chapters@set[2]{\@namedef{chapters@#1}{#2}}
\chapters@set{20261125T001000}{1}% zeros-to-ramp-sums
\chapters@set{20261125T001100}{2}% inert-primes-and-estimate-b
\chapters@set{20261125T005000}{3}% galois-fourier-and-averaging
\chapters@set{20261125T005500}{4}% columns-from-a-sphere
\chapters@set{20261125T000000}{5}% slope-method
\chapters@set{20261126T030100}{6}% a-twist-for-each-prime
\chapters@set{20261125T003000}{7}% the-chain-and-the-square-root
\chapters@set{20261125T010000}{8}% rank-lemma-and-zero-estimate
\chapters@set{20261125T010500}{9}% poincare-series-and-near-orthogonality
\chapters@set{20261124T223000}{10}% elliptic-curves-for-the-proof
\chapters@set{20261125T009500}{11}% fock-model-and-theta-intertwiners
\chapters@set{20261124T223100}{12}% complex-multiplication-and-inert-reduction
\chapters@set{20261125T008000}{13}% one-cm-determinant
\chapters@set{20261124T214500}{14}% supersingular-theta-divisibility
\chapters@set{20261124T214600}{15}% the-isotropic-plane
\chapters@set{20261125T013700}{16}% pairings-on-a-formal-group-of-height-two
\chapters@set{20261125T013500}{17}% hecke-nondegeneracy-on-the-tree
\chapters@set{20261126T030200}{18}% the-bound-is-attained
\chapters@set{20261125T009000}{19}% kronecker-limit-and-chowla-selberg
\chapters@set{20261125T002000}{20}% definite-theta-lifts
\chapters@set{20261125T015500}{21}% siegel-weil-in-the-smallest-case
\chapters@set{20261125T002100}{22}% definite-theta-lifts-and-norms
\chapters@set{20261125T016500}{23}% hecke-operators-and-theta-lifts
\chapters@set{20261125T016700}{24}% divisibility-along-the-chain
\chapters@set{20261125T004000}{25}% denominators-and-sha
\chapters@set{20261125T004100}{26}% tate-twists-and-sha
\chapters@set{20261126T020430}{27}% the-minus-part-by-gauss-sums
\chapters@set{20261125T014100}{28}% congruence-modules
\chapters@set{20261125T014500}{29}% eisenstein-congruences-and-selmer-classes
\chapters@set{20261125T014700}{30}% pseudocharacters-and-the-length-bound
\chapters@set{20261125T011000}{31}% galois-representations-black-box
\chapters@set{20261126T030300}{32}% the-data
\chapters@set{20261126T020210}{33}% bessel-poincare-series-and-their-gram-operators
\chapters@set{20261126T030400}{34}% the-seesaw-identity
\chapters@set{20261126T030500}{35}% divisibility-of-the-composite-lift
\chapters@set{20261126T030600}{36}% theorem-a-short-forms
\chapters@set{20261126T030700}{37}% the-determinant-of-a-chain
\chapters@set{20261126T030800}{38}% the-cost-of-one-lift
\chapters@set{20261126T030900}{39}% denominators-of-projections-to-lifts
\chapters@set{20261126T031000}{40}% twists-with-few-congruences
\chapters@set{20261126T031100}{41}% theorem-c-and-the-main-inequality
\providecommand\chapnumber[1]{%
  \@ifundefined{chapters@#1}%
    {\PackageError{chapters}{Sheet #1 is not listed in the index}%
      {List it in 20261125T012000--grh-workbook__guide.org\MessageBreak
       and run scripts/make-chapters.py.}}%
    {\edef\thechap{\@nameuse{chapters@#1}}}}
\makeatother
\chapnumber{<sheet ID>}
\makeatletter
\providecommand\chapters@set[2]{\@namedef{chapters@#1}{#2}}
\chapters@set{20261125T001000}{1}% zeros-to-ramp-sums
\chapters@set{20261125T001100}{2}% inert-primes-and-estimate-b
\chapters@set{20261125T005000}{3}% galois-fourier-and-averaging
\chapters@set{20261125T005500}{4}% columns-from-a-sphere
\chapters@set{20261125T000000}{5}% slope-method
\chapters@set{20261126T030100}{6}% a-twist-for-each-prime
\chapters@set{20261125T003000}{7}% the-chain-and-the-square-root
\chapters@set{20261125T010000}{8}% rank-lemma-and-zero-estimate
\chapters@set{20261125T010500}{9}% poincare-series-and-near-orthogonality
\chapters@set{20261124T223000}{10}% elliptic-curves-for-the-proof
\chapters@set{20261125T009500}{11}% fock-model-and-theta-intertwiners
\chapters@set{20261124T223100}{12}% complex-multiplication-and-inert-reduction
\chapters@set{20261125T008000}{13}% one-cm-determinant
\chapters@set{20261124T214500}{14}% supersingular-theta-divisibility
\chapters@set{20261124T214600}{15}% the-isotropic-plane
\chapters@set{20261125T013700}{16}% pairings-on-a-formal-group-of-height-two
\chapters@set{20261125T013500}{17}% hecke-nondegeneracy-on-the-tree
\chapters@set{20261126T030200}{18}% the-bound-is-attained
\chapters@set{20261125T009000}{19}% kronecker-limit-and-chowla-selberg
\chapters@set{20261125T002000}{20}% definite-theta-lifts
\chapters@set{20261125T015500}{21}% siegel-weil-in-the-smallest-case
\chapters@set{20261125T002100}{22}% definite-theta-lifts-and-norms
\chapters@set{20261125T016500}{23}% hecke-operators-and-theta-lifts
\chapters@set{20261125T016700}{24}% divisibility-along-the-chain
\chapters@set{20261125T004000}{25}% denominators-and-sha
\chapters@set{20261125T004100}{26}% tate-twists-and-sha
\chapters@set{20261126T020430}{27}% the-minus-part-by-gauss-sums
\chapters@set{20261125T014100}{28}% congruence-modules
\chapters@set{20261125T014500}{29}% eisenstein-congruences-and-selmer-classes
\chapters@set{20261125T014700}{30}% pseudocharacters-and-the-length-bound
\chapters@set{20261125T011000}{31}% galois-representations-black-box
\chapters@set{20261126T030300}{32}% the-data
\chapters@set{20261126T020210}{33}% bessel-poincare-series-and-their-gram-operators
\chapters@set{20261126T030400}{34}% the-seesaw-identity
\chapters@set{20261126T030500}{35}% divisibility-of-the-composite-lift
\chapters@set{20261126T030600}{36}% theorem-a-short-forms
\chapters@set{20261126T030700}{37}% the-determinant-of-a-chain
\chapters@set{20261126T030800}{38}% the-cost-of-one-lift
\chapters@set{20261126T030900}{39}% denominators-of-projections-to-lifts
\chapters@set{20261126T031000}{40}% twists-with-few-congruences
\chapters@set{20261126T031100}{41}% theorem-c-and-the-main-inequality
\providecommand\chapnumber[1]{%
  \@ifundefined{chapters@#1}%
    {\PackageError{chapters}{Sheet #1 is not listed in the index}%
      {List it in 20261125T012000--grh-workbook__guide.org\MessageBreak
       and run scripts/make-chapters.py.}}%
    {\edef\thechap{\@nameuse{chapters@#1}}}}
\makeatother
\chapnumber{20261126T081400}
\renewcommand{\thesection}{\thechap.\arabic{section}}

\theoremstyle{definition}
\newtheorem{theorem}{Theorem}
\renewcommand{\thetheorem}{\thechap.\arabic{theorem}}
\newtheorem{problem}{Problem}[section]
\newtheorem*{leanremark}{In the Lean formalization}

\newcommand{\C}{\mathbb C}
\newcommand{\Z}{\mathbb Z}
\newcommand{\cc}{\mathsf c}
\newcommand{\ee}{\mathsf e}
\newcommand{\ranka}{\mathsf a}
\newcommand{\NN}{\mathsf N}
\newcommand{\kk}{\mathsf k}
\newcommand{\ch}{\mathrm{ch}}
\newcommand{\Ch}{\mathrm{Ch}}
\newcommand{\tr}{\operatorname{tr}}
\newcommand{\diag}{\operatorname{diag}}
\newcommand{\op}{\mathrm{op}}
\newcommand{\tl}[1]{\texttt{#1}}
\setlength{\emergencystretch}{2em}

\title{Black box: discrete series and the comparison of traces}
\author{}
\date{November 26, 2026}

\begin{document}
\maketitle
\refstepcounter{chap}\label{chap}

Two steps of the proof use the same formula for a matrix coefficient at the real places. To form the full Bessel functions of Theorem~\ref{m1-thm:bessel-system}, the proof shifts the coefficient, read in the sign frame, by a chosen position. At each subgroup element it multiplies the shifted coefficient by the inverse subgroup action on a fixed unit vector, including the inverse Bessel character. It averages over the Bessel subgroup and scales the result. The proof of Theorem~\ref{c2-thm:equal-traces} uses conjugacy-class integrals of the corresponding kernels to compare Hecke traces for spaces with different real signatures. Here we state the inputs for both steps and work out the coefficient and its integrals for two-by-two matrices.

\paragraph{Prerequisites.}
Schur orthogonality for compact groups and Weyl's character formula for \(U(n)\).
Problem~\ref{ps-pr:formal-degree} gives the formal degree on the upper half-plane with hyperbolic area. For the comparison we use the framed spaces, standard lattice stabilizers and principal forms from Sections~\ref{ob-sec:spaces} and~\ref{ob-sec:forms}. In this chapter \(K\) denotes a subgroup. We use inner products conjugate-linear in the first argument, as in the definition of the Petersson product in Chapter~\ref{ob-chap}.

\section{One real group}\label{sec:real}

In this chapter let \(n\ge 1\) and let \(\varepsilon=(\varepsilon_i)_{i=1}^n\) have \(p\) plus and \(q\) minus signs. Put \(J=\diag(\varepsilon)\), \(G=U(J)\) and \(K=G\cap U(n)\cong U(p)\times U(q)\).
Let \(l\) have distinct integer entries congruent to \(n-1\) modulo \(2\), and put \(\lambda=l/2\). The positive roots are the coordinate differences \(e_i-e_j\) for \(l_i>l_j\), with \(e_i\) the \(i\)-th coordinate vector. Write \(\rho\) for their half sum and \(\rho_c\) for the half sum over \(\varepsilon_i=\varepsilon_j\). For a nonnegative integer \(h\), we say that \(l\) has gap at least \(h\) if \(|\lambda_i-\lambda_j|\ge h\) for every opposite-sign pair.

\begin{theorem}[Discrete series]\label{thm:ds}
There is an irreducible unitary representation \(\pi_l\) of \(G\) with square-integrable matrix coefficients and Harish-Chandra parameter \(\lambda\). In the restriction to \(K\), the irreducible constituent of highest weight \(\Lambda=\lambda+\rho-2\rho_c\) occurs once and is called the lowest \(K\)-type. The other constituents have highest weights obtained from \(\Lambda\) by adding positive roots. For a unit highest-weight vector \(u\) in that constituent, put \(\cc_l(g)=\langle u,\pi_l(g)u\rangle\), so \(\cc_l(1)=1\).
\end{theorem}

\noindent Existence is classical \cite[Theorem~16]{HC}. For the \(K\)-types see \cite[Theorem~(8.5)]{AS}, \cite[\S2.2]{BV}. For an account of the parametrization see Knapp \cite{Knapp}. The proof uses the integral construction in Theorem~\ref{thm:fj} for parameters of gap at least \(500n^3+10n^5\) in the trace comparison.

Write \(\ch_l\) for the character of the irreducible \(U(n)\)-representation of highest weight \(\lambda-\rho\), with the order on coordinates determined by \(l\). In the compact-dual Haar normalization, \(dg\) on \(G\) has the same density at the identity as probability Haar measure on \(U(n)\) when the tangent vectors are identified by \(X\mapsto JX\). Write \(dk\) for probability Haar measure on \(K\).

\begin{theorem}[Schur orthogonality and formal degree]\label{thm:schur}\label{thm:degree}
With \(d_l=(\int_G|\cc_l(g)|^2\,dg)^{-1}\), for \(x,y\in G\) one has
\[
 d_l\int_G \cc_l(xgy)\overline{\cc_l(g)}\,dg=\cc_l(x)\cc_l(y),
 \qquad
 d_l=\ch_l(1)=
 \frac{\prod_{l_i>l_j}(\lambda_i-\lambda_j)}{\prod_{a=0}^{n-1}a!}.
\]
\end{theorem}

\noindent Both identities are classical. The degree uses the compact-dual Haar normalization \cite[Theorem~(8.1)]{AS}. For the coefficient defined by the integral below, the proof establishes the Schur identity for gap at least \(500n^3\) and the degree formula for gap at least \(500n^3+10n^5\).

\begin{samepage}
\begin{theorem}[Flensted-Jensen's formula]\label{thm:fj}
For an integer weight \(\mathbf w=(w_i)\) and an invertible matrix \(g\), let \(m_{\mathbf w}(g)\) be the product over integers \(s\) of the determinant of the submatrix on rows and columns \(\{i:w_i\ge s\}\), with each factor for \(s\le0\) divided by \(\det g\). The empty minor is \(1\) and only finitely many factors differ from \(1\). For \(g\in G\), keep the two same-sign diagonal blocks of \(g\) and set the other blocks to zero to obtain an invertible matrix \(D_g\). With \(z=D_g^{-1}g\), the integral formula is
\[
 \cc_l(g)=
 \frac{\displaystyle\int_K
   \frac{m_\Lambda(D_g k)\overline{m_\Lambda(k)}}{m_{-\lambda-\rho}(k^*zk)}
   \,dk}
 {\displaystyle\int_K |m_\Lambda(k)|^2\,dk},
\]
where \(k^*\) is the conjugate transpose.
\end{theorem}

\noindent The proof defines the coefficient by the displayed integral. Its identification with the lowest-type coefficient is classical, by the construction of Flensted-Jensen and Sun \cite{FJ,Sun}.
\end{samepage}

\begin{theorem}[Decay]\label{thm:decay}
If \(l\) has gap at least \(h\), then for each integer \(0\le c\le h\) and each \(g\in G\),
\[
 |\cc_l(g)|\le (2/\|g\|_{\op})^c ,
\]
with the Euclidean operator norm. In particular \(|\cc_l(g)|\le 1\).
\end{theorem}

\noindent The estimate with constant \(2\) is read from the proof. For asymptotic decay of discrete-series matrix coefficients see Mili\v{c}i\'c \cite{Mil}.

\begin{theorem}[Character on the diagonal torus]\label{thm:character}
Let \(x=\diag(t_1,\ldots,t_n)\) have distinct entries \(t_i\) of modulus one. In this chapter set \(\kk_i=(l_i+n-1)/2\). Write \(A_K(\kk)(x)\) for the alternating sum of the monomial \(\prod_i t_i^{\kk_i}\) over permutations preserving the signs. The character \(\Ch_l\) of \(\pi_l\) on the regular diagonal torus satisfies
\[
 \ch_l(1)\int_G\cc_l(g^{-1}xg)\,dg
   = \Ch_l(x)
   = (-1)^{pq}\frac{A_K(\kk)(x)}{\prod_{l_i>l_j}(t_i-t_j)} .
\]
Define the kernel \(\ee_l=d_l\overline{\cc_l}\). For gap at least \(500n^3+10n^5\) the proof gives \(d_l=\ch_l(1)\) for the integral coefficient. With \(x_\sigma=\diag(t_{\sigma(1)},\dots,t_{\sigma(n)})\) for each of the \(n!\) permutations \(\sigma\), the proof gives
\[
 (-1)^{pq}\sum_{\sigma\in S_n}\int_G
  \ee_l(g^{-1}x_\sigma g)\,dg
   =p!\,q!\,\overline{\ch_l(x)}.
\]
The classical single orbital integral converges absolutely. The proof gives absolute convergence for the integral coefficient at the gap of the sum.
\end{theorem}

\noindent The first identity uses the classical character and orbital-integral formulas, with the compact centralizer given mass one \cite[Theorem~(8.1)]{AS}, \cite[formula~(1)]{Art}. The sum over permutations is read from the proof.

The next problems check the degree in the compact case and then find the coefficient, measure and degree for \(U(1,1)\). The smallest weights give rank-one models of the large-gap identities.

\begin{problem}\label{pr:compact}
For an integer \(j\ge 0\), the normalized coefficient of the vector \(e_1^{\otimes j}\) in \(\mathrm{Sym}^j\C^2\) is \(g_{11}^j\), with \(e_1=(1,0)\) and the inner product induced from tensor powers of \(\C^2\). For probability Haar measure on \(U(2)\), \(\int_{U(2)} |g_{11}|^{2j}\,dg=1/(j+1)\), equal to \(1/2\) at \(j=1\).
\end{problem}

In the rank-one problems \(\ranka\ge1\) and \(\NN\ge2\) are integers, and \(K\) is the diagonal torus of \(U(1,1)\). Use the same diagonal torus in \(U(2)\).

\begin{problem}\label{pr:weight}
For \(G=U(1,1)\) with signs \((+,-)\), the parameter \(l=(-2\ranka-1,-2\ranka-2\NN+1)\) gives the lowest weight \((-\ranka,-\ranka-\NN)\) and maximal gap \(\NN-1\).
\end{problem}

\begin{problem}\label{pr:coefficient}
For the parameter in Problem~\ref{pr:weight}, Theorem~\ref{thm:fj} gives \(\cc_l(g)=\det(g)^{-\ranka}g_{22}^{-\NN}\). The coefficient transforms on the left and the right under \(K\) by the character \( (u_1,u_2)\mapsto u_1^{-\ranka}u_2^{-\ranka-\NN}\).
\end{problem}

In the chart of complex lines \([z:1]\), \(U(2)/K\) is the sphere and \(U(1,1)/K\) is the disk \(|z|<1\), by the action on the line \([0:1]\). The two invariant measures have densities proportional to \((1+|z|^2)^{-2}\) and \((1-|z|^2)^{-2}\), respectively, against Euclidean area. For the compact-dual normalization the two quotient densities agree at \(z=0\). Write
\[
 a_t=\begin{pmatrix}\cosh t&\sinh t\\\sinh t&\cosh t\end{pmatrix}\qquad(t\ge0).
\]

\begin{problem}\label{pr:radial}
In the compact chart take \(|z|=\tan t\), \(0\le t<\pi/2\). Probability measure has radial density \(\sin(2t)\,dt\). In the disk take \(|z|=\tanh t\), the radius of \(a_t[0:1]\). The compact-dual measure has radial density \(\sinh(2t)\,dt\).
\end{problem}

\begin{problem}\label{pr:degree}
The squared mass of the coefficient in Problem~\ref{pr:coefficient} is
\[
 \int_0^\infty \cosh(t)^{-2\NN}\sinh(2t)\,dt=\frac{1}{\NN-1}.
\]
At \(\NN=3\) the mass is \(1/2\). The reciprocal of the mass is \((l_1-l_2)/2=\ch_l(1)\).
\end{problem}

\begin{problem}\label{pr:half-plane}
The pullback of half-plane area \(dx\,dy/y^2\) under \(z\mapsto \mathrm{i}(1+z)/(1-z)\) is \(4\pi\) times the compact-dual quotient measure on the disk. The rescaled degree is \((\NN-1)/(4\pi)\). For even \(\NN\ge4\) the result agrees with Problem~\ref{ps-pr:formal-degree} at weight \(\NN\).
\end{problem}

\begin{problem}\label{pr:rank-decay}
For \(g=k a_t k'\) with \(k,k'\in K\), the coefficient in Problem~\ref{pr:coefficient} has modulus \(\cosh(t)^{-\NN}\), and \(\|g\|_{\op}=e^t\). In particular \(|\cc_l(g)|\le (2/\|g\|_{\op})^{\NN}\).
\end{problem}

The last problems in the real group section check the sign in the torus formula and the character recovered by conjugation averaging.

\begin{problem}\label{pr:sign}
For signs \((+,-)\) and \(\kk_1>\kk_2\), write
\(\Ch(t_1,t_2)=t_1^{\kk_1}t_2^{\kk_2}/(t_2-t_1)\) for the character in Theorem~\ref{thm:character} at distinct unit entries. Then
\[
 -\bigl(\Ch(t_1,t_2)+\Ch(t_2,t_1)\bigr)
  =\sum_{i=0}^{\kk_1-\kk_2-1} t_1^{\kk_1-1-i}t_2^{\kk_2+i}.
\]
The sum is the \(U(2)\)-character of highest weight \((\kk_1-1,\kk_2)\), and is \(t_1^2+t_1t_2+t_2^2\) at \((\kk_1,\kk_2)=(3,0)\). The local sign is \((-1)^{pq}=-1\). With one fixed signature at an even number \(f\) of infinite places, the product of the local signs is \(1\).
\end{problem}

\begin{problem}\label{pr:finite-orbital}
Let \(\pi\) be an irreducible unitary representation of a finite group \(G_0\), of dimension \(d_\pi\), and \(u\) a unit vector. For \(\cc(g)=\langle u,\pi(g)u\rangle\), one has
\[
 \frac{d_\pi}{\#G_0}\sum_{g\in G_0}\cc(g^{-1}xg)=\tr\pi(x)\quad(x\in G_0).
\]
\end{problem}

\begin{problem}\label{pr:s3}
For the action of \(S_3\) permuting the coordinates on the sum-zero subspace of \(\C^3\), the averages on the left of Problem~\ref{pr:finite-orbital} are \(2,-1,0\) at the identity, a three-cycle and a transposition, respectively.
\end{problem}

The comparison section gives the vanishing and trace identities for framed spaces.

\section{The comparison}\label{sec:comparison}

Let \(E\) be a CM field with a fixed CM type and \(V\) a framed space of rank \(n\ge2\) as in Section~\ref{ob-sec:spaces}. The frames identify the finite adelic unitary groups of all such spaces with the group \(G_f\) for the identity form. For each infinite place \(\tau\), choose a doubled parameter \(l_\tau\) for the signs of \(V\) of gap at least \(500n^3+10n^5\). Use the product of the corresponding kernels \(\ee_{l_\tau}\) on the real factors.

\begin{theorem}[Selberg principle]\label{thm:selberg}
Let \(\phi\) be a locally constant compactly supported function on \(G_f\). For a rational unitary element \(\gamma\) of \(V\), integrate the product of the real kernels and \(\phi\), evaluated at the real and finite factors of \(g^{-1}\gamma g\), over the left quotient of the adelic group by the rational centralizer of \(\gamma\). Use the quotient of product Haar measure by counting measure on the rational centralizer. At infinity use the compact-dual measures, and on \(G_f\) give the standard lattice stabilizer mass one. The orbital integral is zero if the characteristic polynomial over \(E\) is irreducible and has a root off the unit circle under the embedding selected by the CM type at some infinite place.
\end{theorem}

\noindent The adelic statement is read from the proof. The local vanishing for discrete-series coefficients on regular nonelliptic classes is classical \cite[formula~(1)]{Art}.

\begin{theorem}[Comparison of traces]\label{thm:traces}
With \(E\), \(V\) and \(l_\tau\) as in the setup for this section, suppose that
\begin{enumerate}
\item \(E\) contains the \(5040\)th roots of unity and has an even number of infinite places.
\item \(V'\) is a framed space of the same rank over \(E\), and the parameters \(l_\tau\) meet the same gap condition for the signs of \(V'\).
\item \(K_f\) is a subgroup of the integral elements of \(G_f\) containing a neighbourhood of the identity, and \(a\in\C[G_f]\) preserves \(K_f\)-fixed vectors in every representation of \(G_f\).
\item The function \(\phi=a\star 1_{K_f}\) is cuspidal at a split place \(w\). Here \(\star 1_{K_f}\) sends each \(b\in G_f\) to \(1_{bK_f}\) and extends linearly. For cuspidality, let \(P\) be the coordinate projection onto any nonempty proper coordinate subset and let \(\bar b\in\mathrm{GL}_n(k_w)\), where \(k_w\) is the residue field. For each matrix \(A\in M_n(k_w)\) take any integral invertible lift \(u_A\) of \((1+PA(1-P))\bar b\), viewed in \(G_f\) with inverse conjugate transpose at the conjugate place and identity components elsewhere. For every choice of lifts and every \(g\in G_f\), the sum of \(\phi(gu_A)\) over \(A\) is zero.
\end{enumerate}
Then the operators induced by \(a\) through right translation have equal traces on the spaces of forms for \(V\) and \(V'\). The spaces are defined as for principal forms in Section~\ref{ob-sec:forms}, using level \(K_f\) and the real kernels of the parameters \(l_\tau\) for the respective signs.
\end{theorem}

\noindent The method is classical \cite{AC,FK,Lab}; for the use of a cuspidal component see Harris--Labesse \cite[Proposition~1.10.5]{HL}. The statement with the discrete-series kernels and the split cuspidality condition above is read from the proof. Chapter~\ref{c2-chap} checks the comparison hypotheses for the good forms.

\begin{theorem}[Landherr's theorem]\label{thm:landherr}
Hermitian spaces for a quadratic extension of number fields that are isometric at each completion of the base field are isometric over the extension.
\end{theorem}

\noindent The theorem is classical \cite[\S5.5]{Kne}. The proof establishes the case of framed spaces with the same numbered signs over a CM field containing the \(5040\)th roots of unity.

\begin{theorem}[Bessel models]
Use the Bessel subgroup and character of Section~\ref{bp-sec:subgroup} at a real place. A Bessel functional is a continuous linear form on the completed tensor product of two smooth representations, equivariant for the generic Bessel character. The subgroup acts on the smaller representation through the quotient by its unipotent radical. For an irreducible pair of smooth admissible Fr\'echet representations of moderate growth with finite-length modules of \(K\)-finite vectors, the space of Bessel functionals has dimension at most one. For fixed tempered Langlands parameters for the two groups, exactly one pair across the relevant pure inner forms has a nonzero Bessel functional.
\end{theorem}

\noindent The statements are context from the literature \cite[Theorem~A]{JSZ}, \cite{BP,Xue}. In codimension one for discrete series see also He \cite{He}. In Chapter~\ref{b1-chap}, Theorem~\ref{b1-thm:period} gives the lower bound for the Bessel period used in normalizing the real Bessel functions.

\medskip
The problems check the normalization, degree, decay and character sign on the two-by-two groups and conjugation averaging on a finite group. We took as given the discrete-series parametrization and lowest type, Schur orthogonality and formal degree, the principal-minor formula, the general decay estimate, the character formulas, Selberg vanishing, trace comparison, Landherr's theorem and the statements about Bessel models.

\begin{leanremark}
\raggedright
In \cite{GL}, the parameter is \tl{HCParam}. The condition \tl{HCParam.HasLargeGap} asks for gap at least \(500n^3+A\) with \(A\) a nonnegative integer. The integral coefficient is \tl{HCParam.dsCoeff}. For Schur orthogonality and the degree formula see \tl{dsCoeff\_schur\_orthogonality} and \tl{Orbital.formalDegree\_eq\_weylDimensionPoly}. For the decay and the sum over permutations see \tl{DiscreteSeries.norm\_dsCoeff\_le\_pow} and \tl{Orbital.neg\_one\_pow\_mul\_sum\_integral\_dsKernel}. Vanishing, comparison and the framed case of Landherr are \tl{Space.orbitalIntegral\_eq\_zero}, \tl{trace\_eq\_of\_isCuspidalAt} and \tl{Space.landherr}.
\end{leanremark}

\section*{Hints}
\begin{description}
\item[\ref{pr:coefficient}] First specialize the product of minors to diagonal matrices and matrices with diagonal entries \(1\).
\item[\ref{pr:radial}] Use Euclidean polar area, and fix the sphere constant by total mass one.
\item[\ref{pr:finite-orbital}] Schur's lemma.
\end{description}

{\small
\begin{thebibliography}{99}
\bibitem[AC]{AC} J.~Arthur and L.~Clozel, \emph{Simple Algebras, Base Change, and the Advanced Theory of the Trace Formula}, Ann. of Math. Stud. 120, Princeton University Press, 1989.
\bibitem[Ar]{Art} J.~Arthur, \emph{The characters of discrete series as orbital integrals}, Invent. Math. \textbf{32} (1976), 205--261.
\bibitem[AS]{AS} M.~F.~Atiyah and W.~Schmid, \emph{A geometric construction of the discrete series for semisimple Lie groups}, Invent. Math. \textbf{42} (1977), 1--62.
\bibitem[BV]{BV} V.~Baldoni and M.~Vergne, \emph{Discrete series representations and \(K\) multiplicities for \(U(p,q)\). User's guide}, arXiv:1008.5360, 2010.
\bibitem[BP]{BP} R.~Beuzart-Plessis, \emph{A local trace formula for the Gan--Gross--Prasad conjecture for unitary groups: the Archimedean case}, Ast\'erisque \textbf{418} (2020).
\bibitem[FK]{FK} Y.~Z.~Flicker and D.~A.~Kazhdan, \emph{A simple trace formula}, J. Analyse Math. \textbf{50} (1988), 189--200.
\bibitem[FJ]{FJ} M.~Flensted-Jensen, \emph{Discrete series for semisimple symmetric spaces}, Ann. of Math. \textbf{111} (1980), 253--311.
\bibitem[HC]{HC} Harish-Chandra, \emph{Discrete series for semisimple Lie groups. II. Explicit determination of the characters}, Acta Math. \textbf{116} (1966), 1--111.
\bibitem[He]{He} H.~He, \emph{On the Gan--Gross--Prasad conjectures for \(U(p,q)\)}, Invent. Math. \textbf{209} (2017), 837--884.
\bibitem[HL]{HL} M.~Harris and J.-P.~Labesse, \emph{Conditional base change for unitary groups}, Asian J. Math. \textbf{8} (2004), 653--684.
\bibitem[JSZ]{JSZ} D.~Jiang, B.~Sun and C.-B.~Zhu, \emph{Uniqueness of Bessel models: the archimedean case}, Geom. Funct. Anal. \textbf{20} (2010), 690--709.
\bibitem[Knapp]{Knapp} A.~W.~Knapp, \emph{Representation Theory of Semisimple Groups. An Overview Based on Examples}, Princeton University Press, 1986.
\bibitem[Kn]{Kne} M.~Kneser, \emph{Lectures on Galois Cohomology of Classical Groups}, Tata Inst. Lectures 47, Bombay, 1969.
\bibitem[Lab]{Lab} J.-P.~Labesse, \emph{Changement de base CM et s\'eries discr\`etes}, in \emph{On the Stabilization of the Trace Formula}, International Press, 2011, 429--470.
\bibitem[Mi]{Mil} D.~Mili\v{c}i\'c, \emph{Asymptotic behaviour of matrix coefficients of the discrete series}, Duke Math. J. \textbf{44} (1977), 59--88.
\bibitem[GL]{GL} L.~Sernau, \emph{GRH for Dirichlet \(L\)-functions, in Lean~4}, a shorter rewriting of the Lean proof, text by Claude, \url{https://github.com/lsernau/grh-lean}, commit \texttt{72fead1}, readable edition \texttt{grh/}.
\bibitem[Sun]{Sun} B.~Sun, \emph{Matrix coefficients of cohomologically induced representations}, Compos. Math. \textbf{143} (2007), 201--221.
\bibitem[Xue]{Xue} H.~Xue, \emph{Bessel models for real unitary groups: the tempered case}, Duke Math. J. \textbf{172} (2023), 995--1031.
\end{thebibliography}
}

\end{document}
